\chapter{Quantum Fluctuations and the Origin of Structure}

\lectureoverview{The final lecture asks how a universe made smooth by inflation can nevertheless acquire the tiny irregularities from which galaxies grow. The answer is built in stages. Universal constants first define the Planck scales and expose the enormous gap between familiar particle physics and quantum gravity. The microwave sky then supplies the target: temperature and density variations of order \(10^{-5}\), spread over many scales. Zero-point motion supplies unavoidable fluctuations, damping explains how an oscillation can freeze, and the expanding-universe wave equation turns every Fourier mode of the inflaton into just such a damped oscillator. Modes freeze when their physical wavelength reaches the Hubble scale. When inflation ends, the frozen variation makes neighboring regions reheat at slightly different times, converting microscopic quantum noise into the primordial density pattern.}

\section{Which Constants Are Universal?}

We possess an extraordinarily accurate quantum theory of elementary particles. Its parameters include quark and lepton masses, coupling constants, mixing data, charges, and the list of particles themselves. Given unlimited calculating power, the theory would account not only for accelerator experiments but also for atoms, molecules, and chemistry; in principle the consequences continue into biology and the physical processes underlying thought. That does not make such systems easy to calculate. It says that the Standard Model is intended as a microscopic foundation for ordinary matter. Gravity is described separately, at the classical level, by general relativity. Both theories are powerful; neither makes every number appearing in low-energy physics look inevitable.

A proton-to-electron mass ratio is a perfectly good physical number, but it is tied to two particular particles. Replacing the proton by a neutron or the electron by a muon produces a different ratio. Such parameters may encode deep physics that we do not yet know, but they do not advertise their universality. Three dimensional constants are different:
\[
c,\qquad \hbar,\qquad G.
\]
Their defining laws repeatedly contain the word \emph{all}. All matter and information respect the limiting speed \(c\). All quantum systems obey uncertainty relations governed by \(\hbar\). All forms of energy gravitate with the same Newton constant \(G\). That breadth is what makes these constants fundamental in the present argument.

Planck recognized the importance of this combination before the full meaning of quantum mechanics was known. His constant first entered the description of blackbody radiation. Even then, the dimensional fact was visible: from \(c\), \(\hbar\), and \(G\) one can construct units of length, time, and mass without referring to any particular atom or particle.

\section{Building the Planck Scales}

Let \(L\), \(T\), and \(M\) denote the dimensions of length, time, and mass. The speed of light has dimensions
\[
[c]=LT^{-1}.
\]
The uncertainty relation \(\Delta x\,\Delta p\gtrsim\hbar\), together with \(p=mv\), gives
\[
[\hbar]=ML^2T^{-1}.
\]
Finally, Newton's equations
\[
F=\frac{Gm_1m_2}{r^2},
\qquad
F=ma
\]
give
\[
[G]=L^3M^{-1}T^{-2}.
\]
This derivation matters because dimensional analysis is not guesswork: each exponent is fixed by requiring unwanted dimensions to cancel.

Seek a length of the form
\begin{equation}
\ell=c^\alpha\hbar^\beta G^\gamma.
\label{eq:c10-length-ansatz}
\end{equation}
The mass exponent must vanish, so
\[
\beta-\gamma=0.
\]
The time exponent must vanish, and the length exponent must equal one:
\[
-\alpha-\beta-2\gamma=0,
\qquad
\alpha+2\beta+3\gamma=1.
\]
Using \(\gamma=\beta\), these become
\[
\alpha+3\beta=0,
\qquad
\alpha+5\beta=1.
\]
Their difference gives \(2\beta=1\), hence
\[
\beta=\gamma=\frac12,
\qquad
\alpha=-\frac32.
\]
Therefore the unique length made from the three universal constants is
\begin{equation}
\ell_{\mathrm P}
=\sqrt{\frac{\hbar G}{c^3}}
\approx 1.6\times10^{-35}\,\mathrm m.
\label{eq:c10-planck-length}
\end{equation}
The lecture keeps only the order of magnitude, \(10^{-35}\,\mathrm m\), but the dimensional derivation fixes the combination itself.

The associated time is the light-crossing time:
\begin{equation}
t_{\mathrm P}
=\frac{\ell_{\mathrm P}}{c}
=\sqrt{\frac{\hbar G}{c^5}}
\approx 5.4\times10^{-44}\,\mathrm s.
\label{eq:c10-planck-time}
\end{equation}
At order-of-magnitude accuracy this is \(10^{-43}\,\mathrm s\). It is not an independent construction once the Planck length and the universal speed are known.

For the mass, use the dimensional relation \(Et\sim\hbar\) and \(E=mc^2\). Taking \(t=t_{\mathrm P}\) gives
\begin{equation}
m_{\mathrm P}
=\frac{\hbar}{c^2t_{\mathrm P}}
=\sqrt{\frac{\hbar c}{G}}
\approx 2.2\times10^{-8}\,\mathrm{kg}.
\label{eq:c10-planck-mass}
\end{equation}
This is surprisingly macroscopic by particle standards, roughly the mass of a fine dust mote. Its rest energy is of order \(10^9\,\mathrm J\), comparable in scale to the chemical energy in a tank of fuel. A car can travel hundreds of miles on only a few tanks because combustion releases merely a minute fraction of the fuel's rest energy. Converting one Planck mass completely into energy would be an entirely different affair. A tiny mass can represent an enormous energy when multiplied by \(c^2\).

The contrast with accelerator physics is severe. A distance of order \(10^{-18}\,\mathrm m\), characteristic of the shortest scales probed in high-energy collisions at the time of the lecture, is still about \(10^{17}\) Planck lengths. The unknown territory between the electroweak and Planck scales is therefore vast. Nuclear physics offers the right analogy. A nucleus looks elementary from far away, but at shorter distances it resolves into protons and neutrons, and these in turn resolve into quarks and gluons. Low-energy particles may be simple compared with nuclei yet still be complicated emergent objects relative to whatever lies at the Planck scale. The many parameters of the Standard Model are less surprising when viewed across that hierarchy, though their observed values and fine tunings remain unexplained.

\begin{classroomqa}
\audiencequestion{Is the Higgs lighter or heavier than expected, and why is its mass puzzling?}
\lecturerresponse{There are two comparisons. Relative to the Planck mass, a \(125\,\mathrm{GeV}\) Higgs is extraordinarily light. In a theory with no protective structure, quantum corrections tend to make a scalar mass sensitive to much higher scales, so obtaining the observed value looks like balancing parameters on a narrow edge. Relative to some simple supersymmetric expectations discussed at the time, the observed Higgs could instead be called somewhat heavy. The apparent contradiction comes from comparing it with two different reference scales; the underlying hierarchy problem remains. \lecturetimestamp{00:26:45--00:30:54}}
\end{classroomqa}

\begin{classroomqa}
\audiencequestion{If the Higgs helps give particles mass, should dark-matter particles have masses near the Planck scale?}
\lecturerresponse{No. A WIMP is a \emph{weakly interacting massive particle}; the phrase points to an interaction scale that might make laboratory detection possible, not to the Planck mass. Under the usual WIMP assumptions, masses can be well above the proton mass and still be fantastically below \(m_{\mathrm P}\). If the interaction with ordinary matter is allowed to be arbitrarily weaker, cosmology permits a much wider mass range, but direct discovery correspondingly becomes harder. The proton and neutron also provide a useful warning: most of their mass is binding energy from the strong interaction, not simply a sum of Higgs-generated quark masses. \lecturetimestamp{00:30:55--00:33:11}}
\end{classroomqa}

\begin{classroomqa}
\audiencequestion{What does a Planck mass signify, and is there a fundamental energy per bit?}
\lecturerresponse{A useful gravitational interpretation is that a black hole cannot be made parametrically smaller than the Planck scale while remaining a conventional semiclassical black hole. Such an object has mass of order \(m_{\mathrm P}\), horizon area of order one Planck area, and entropy of order one unit. In that restricted black-hole sense, one may associate an order-Planck energy with order-one information. This is not a universal claim that every stored bit in every device costs a Planck energy; ordinary particles and ordinary memories can be far lighter. \lecturetimestamp{00:33:11--00:35:10}}
\end{classroomqa}

From here it is convenient to use Planck units,
\begin{equation}
c=\hbar=G=1.
\label{eq:c10-planck-units}
\end{equation}
Lengths, times, masses, and energies can then be quoted as pure numbers measured against equations~\eqref{eq:c10-planck-length}--\eqref{eq:c10-planck-mass}.

\section{The Nearly Uniform Microwave Sky}

The cosmological problem begins with a microwave map of the whole sky. WMAP, the Wilkinson Microwave Anisotropy Probe, measured radiation released near the epoch when the primordial plasma became transparent. The photons have since been redshifted by roughly a factor of one thousand and now form an almost perfect blackbody at about \(2.7\,\mathrm K\). We do not see the universe at one common present time when we look outward. Greater distance means earlier emission, and the microwave background is the distant shell from which photons first streamed freely toward us.

\begin{classroomqa}
\audiencequestion{What exactly does the oval WMAP image represent?}
\lecturerresponse{Every point in the oval represents a direction on the sky, with the sphere flattened by a map projection. The left and right edges join, the upper and lower portions represent opposite galactic hemispheres, and the galactic plane runs across the middle. Foreground microwave emission from our own Galaxy must be modeled and subtracted. What remains is not a map of nearby stars but a directional map of the microwave temperature on the last-scattering shell. \lecturetimestamp{00:37:49--00:40:49}}
\end{classroomqa}

To first approximation the result is isotropic:
\[
T(\hat{\mathbf n})\simeq 2.7\,\mathrm K
\]
for every viewing direction \(\hat{\mathbf n}\). Perfect isotropy around every point would imply perfect homogeneity. But a perfectly homogeneous universe would not spontaneously produce galaxies, clusters, and the caustic-like cosmic web discussed in the previous lecture. Some inhomogeneity must have been present, however small its amplitude.

Write the directional temperature departure as
\[
\delta T(\hat{\mathbf n})=T(\hat{\mathbf n})-\bar T.
\]
The measured fractional variation is of order
\begin{equation}
\frac{\delta T}{T}\sim10^{-5}.
\label{eq:c10-cmb-anisotropy}
\end{equation}
One route to this number was dynamical: evolve today's matter distribution backward and ask what initial contrast is needed at decoupling. The other was direct: amplify the nearly featureless microwave map by five orders of magnitude and measure the mottling itself. The two routes agree at the level needed for this argument.

The mottling occurs on many angular scales. Large patches contain smaller patches, and magnifying a region gives a statistically similar texture. The term \emph{scale invariant} does not mean that every scale is literally identical; it means that the fluctuation power changes only weakly with scale over the relevant range. That near scale invariance is part of the evidence the theory must explain, not merely the overall amplitude.

Why should temperature trace matter? A photon escaping a gravitational potential well loses energy. Radiation arriving from a denser region can therefore be shifted relative to radiation from a less dense region. At the lecture's order-of-magnitude level this motivates
\begin{equation}
\frac{\delta\rho}{\rho}\sim10^{-5}
\label{eq:c10-density-target}
\end{equation}
near decoupling. A precise CMB calculation includes acoustic dynamics, gravitational potentials, and transfer functions, so equation~\eqref{eq:c10-density-target} is a scale-setting shorthand rather than an exact equality between \(\delta T/T\) and \(\delta\rho/\rho\).

\begin{classroomqa}
\audiencequestion{Could the mottling be ordinary microscopic noise generated at last scattering?}
\lecturerresponse{Not on the observed angular scales. The patches corresponded to enormous physical regions by decoupling, not atomic-scale disturbances in the recombining plasma. Tracing them much farther backward, however, makes their physical wavelengths smaller. Inflation supplies the missing magnification: fluctuations that began as quantum, microscopic disturbances can be stretched to astronomical scales. \lecturetimestamp{00:46:42--00:48:02}}
\end{classroomqa}

The mechanism can be understood from three pieces of physics: zero-point motion of a harmonic oscillator, damping of an oscillator, and the wave equation for a scalar field in an expanding universe. We shall assemble them in that order.

\section{Zero-Point Motion of an Oscillator}

For a unit-mass harmonic oscillator,
\begin{equation}
E=\frac12\dot x^{\,2}+\frac12\omega^2x^2,
\qquad
\ddot x+\omega^2x=0.
\label{eq:c10-undamped}
\end{equation}
Kinetic and potential energies continually exchange. Averaged over an oscillation, each carries half the total energy.

Quantum mechanics gives the energy levels
\begin{equation}
E_n=\left(n+\frac12\right)\hbar\omega.
\label{eq:c10-oscillator-levels}
\end{equation}
The ground state cannot have both \(x=0\) and \(\dot x=0\), because that would assign exact position and momentum simultaneously. Its irreducible energy,
\[
E_0=\frac12\hbar\omega,
\]
is the zero-point energy. The word \emph{oscillation} should be understood statistically: a stationary ground state does not trace a classical sinusoidal orbit, but it has nonzero variances in position and momentum.

Half the ground-state energy is in the potential term on average:
\[
\frac12\omega^2\langle x^2\rangle
=\frac14\hbar\omega.
\]
Thus
\begin{equation}
\langle x^2\rangle=\frac{\hbar}{2\omega}
\sim\frac{\hbar}{\omega}.
\label{eq:c10-zero-point-position}
\end{equation}
Ignoring factors of two, a complex notation for the scale of the fluctuation is
\[
x(t)\sim\sqrt{\frac{\hbar}{\omega}}\,e^{i\omega t}.
\]
Differentiation then gives
\begin{equation}
\langle\dot x^{\,2}\rangle\sim\hbar\omega,
\label{eq:c10-zero-point-velocity}
\end{equation}
consistent with equal average kinetic and potential energies. Cooling can remove excitations with \(n>0\), but it cannot remove these ground-state variances.

\section{Damping, Criticality, and Freezing}

Immerse the oscillator in a viscous medium. A spring and mass moving through honey obey
\begin{equation}
\ddot x+\gamma\dot x+\omega^2x=0,
\label{eq:c10-damped}
\end{equation}
where \(\gamma>0\) measures the drag. Looking for \(x=e^{\alpha t}\) yields
\begin{equation}
\alpha^2+\gamma\alpha+\omega^2=0,
\end{equation}
and therefore
\begin{equation}
\alpha_\pm
=\frac{-\gamma\pm\sqrt{\gamma^2-4\omega^2}}{2}.
\label{eq:c10-damped-roots}
\end{equation}
The two roots provide the two independent solutions required by a second-order differential equation.

\begin{classroomqa}
\audiencequestion{Should the drag term in equation~\eqref{eq:c10-damped} carry a minus sign?}
\lecturerresponse{The drag force itself is \(F_{\rm drag}=-\gamma\dot x\). Newton's equation is \(\ddot x=-\gamma\dot x-\omega^2x\). Moving both force terms to the left gives equation~\eqref{eq:c10-damped} with plus signs. The sign in the differential equation and the direction of the physical force are therefore consistent. \lecturetimestamp{00:59:38--01:00:10}}
\end{classroomqa}

If \(\gamma^2<4\omega^2\), the roots are complex. Define
\[
\Omega=\sqrt{\omega^2-\frac{\gamma^2}{4}}.
\]
Then
\begin{equation}
x(t)=e^{-\gamma t/2}
\left(A\cos\Omega t+B\sin\Omega t\right).
\label{eq:c10-underdamped}
\end{equation}
The oscillator still crosses the equilibrium point repeatedly, but its envelope decays. This is the underdamped regime.

If \(\gamma^2>4\omega^2\), both roots in equation~\eqref{eq:c10-damped-roots} are real and negative:
\begin{equation}
x(t)=A e^{\alpha_+t}+B e^{\alpha_-t}.
\label{eq:c10-overdamped}
\end{equation}
There is no sustained oscillation; both components decay. This is the overdamped regime. At
\begin{equation}
\gamma=2\omega
\label{eq:c10-critical}
\end{equation}
the roots coincide, marking critical damping, the crossover between the two behaviors.

One limiting case is especially important. With no restoring force, \(\omega=0\),
\[
\alpha(\alpha+\gamma)=0,
\]
so
\begin{equation}
x(t)=c_0+c_1e^{-\gamma t}.
\label{eq:c10-no-restoring}
\end{equation}
The velocity dies, but the coordinate does not have to return to zero. It freezes at the constant \(c_0\), selected by the initial conditions. A stone pushed through honey with no spring attached comes to rest wherever the drag exhausts its motion.

Now let the frequency vary slowly. Begin with \(\omega\gg\gamma\), so the motion oscillates. Reduce \(\omega(t)\) while keeping the damping scale approximately fixed. The oscillation slows, crosses \(\gamma=2\omega\), becomes overdamped, and approaches the force-free form as \(\omega\to0\). Equivalently, one may raise the damping relative to the restoring force. In the cosmological application the decisive change will be the fall of the effective frequency while the Hubble damping remains nearly constant. The final frozen value cannot be inferred from the crossover alone; it follows from the complete earlier evolution and initial quantum state.

\section{Scalar Waves in an Expanding Universe}

Return to the inflaton \(\phi\). Its potential is assumed to contain a high, broad plateau followed by a steep descent into a lower minimum. On the plateau the important approximation is
\begin{equation}
V'(\phi)\simeq0,
\label{eq:c10-flat-potential}
\end{equation}
not \(V(\phi)\simeq0\). The slope is small while the energy density can remain large enough to drive inflation. The detailed origin of this shape is not supplied; it is a phenomenological requirement of the model.

\begin{figure}[htbp]
\centering
\includegraphics[width=0.88\textwidth]{lecture_10_figure_02.png}
\caption{The inflaton plateau and cliff remain visible at left while the flat-space wave equation is begun at right. The board securely supports the approximate \(V'(\phi)=0\) condition and the start of a second-derivative term, but the lecturer obscures the unfinished equation; equation~\eqref{eq:c10-flat-wave} is reconstructed from the spoken derivation.}
\label{fig:c10-potential-wave}
\lectureframe{01:17:17.865}
\end{figure}

For a nearly free scalar in flat spacetime,
\begin{equation}
\frac{\partial^2\phi}{\partial t^2}
-\nabla^2\phi=0.
\label{eq:c10-flat-wave}
\end{equation}
A nonflat potential would add \(-V'(\phi)\) on the right. For the plateau approximation we omit it. If the field were perfectly homogeneous, the spatial derivative would vanish. In an expanding universe the remaining homogeneous equation would be
\[
\ddot\phi+3H\dot\phi=0,
\]
the Hubble-friction equation of the previous lecture.

\begin{figure}[htbp]
\centering
\includegraphics[width=0.88\textwidth]{lecture_10_figure_03.png}
\caption{The plateau is indicated at left while successive stages of the wave equation appear at right. The upper flat-space operator is readable; the lower line reaches \(+3H\) but is unfinished, so the complete damping term is supplied from the contemporaneous explanation.}
\label{fig:c10-hubble-board}
\lectureframe{01:18:57.554}
\end{figure}

\begin{figure}[htbp]
\centering
\begin{coursefitfigure}
\begin{tikzpicture}[x=1cm,y=0.9cm,>=Latex]
  \draw[->] (-0.2,0) -- (9.3,0) node[right] {\(\phi\)};
  \draw[->] (0,-0.2) -- (0,4.6) node[above] {\(V(\phi)\)};
  \draw[very thick,blue!55!black,smooth]
    plot coordinates {(0.25,3.9) (1.4,3.88) (2.8,3.8) (4.2,3.6) (5.3,3.25) (6.0,2.6) (6.5,1.45) (6.9,0.45) (7.5,0.28) (8.5,0.7)};
  \draw[dashed] (6.45,0) -- (6.45,2.0);
  \draw[->,thick,red!65!black] (2.0,4.25) -- (3.1,3.82);
  \node[anchor=south west] at (0.8,4.0) {slow-roll plateau: \(V'\simeq0\)};
  \node[anchor=west] at (6.55,2.05) {end of inflation};
  \node[anchor=north] at (7.45,0.22) {reheating region};
\end{tikzpicture}
\end{coursefitfigure}
\caption{Clean reconstruction of the potential used throughout the argument. A shallow slope permits slow roll while the high value of \(V\) keeps \(H\) nearly constant; the steep descent ends inflation and releases the stored energy.}
\label{fig:c10-potential-reconstruction}
\end{figure}

For an inhomogeneous field, coordinate distance and physical distance must be separated. Choose dimensionless comoving labels \(\mathbf x\), so an infinitesimal proper distance is
\[
d\boldsymbol\ell=a(t)\,d\mathbf x.
\]
At fixed time,
\[
\nabla_{\rm proper}=\frac{1}{a}\nabla_{\mathbf x},
\qquad
\nabla_{\rm proper}^2=\frac{1}{a^2}\nabla_{\mathbf x}^2.
\]
The scalar equation is therefore
\begin{equation}
\ddot\phi+3H\dot\phi-\frac{1}{a^2}\nabla^2\phi=0,
\label{eq:c10-frw-wave}
\end{equation}
where \(\nabla\) differentiates with respect to the comoving labels. If instead one chooses a dimensionless scale factor, the comoving coordinates carry length; the physics is unchanged.

\begin{classroomqa}
\audiencequestion{Should the spatial term in equation~\eqref{eq:c10-frw-wave} be positive?}
\lecturerresponse{No. The flat-space wave operator already has the relative sign \(\partial_t^2-\nabla^2\). Expansion changes the spatial derivative by the factor \(a^{-2}\), but it does not reverse that sign. After a Fourier substitution the two minus signs combine, and the mode equation acquires the expected positive restoring term \(+k^2\phi_k/a^2\). \lecturetimestamp{01:22:18--01:23:06}}
\end{classroomqa}

\section{Each Fourier Mode Is a Damped Oscillator}

Study one spatial wave at a time:
\begin{equation}
\phi(\mathbf x,t)
=\phi_{\mathbf k}(t)e^{i\mathbf k\cdot\mathbf x}.
\label{eq:c10-mode-ansatz}
\end{equation}
The complex exponential is bookkeeping for sines and cosines. An arbitrary real profile is built by superposing modes with the appropriate reality relation between \(\mathbf k\) and \(-\mathbf k\).

Time derivatives act only on \(\phi_{\mathbf k}(t)\), while
\[
\nabla^2e^{i\mathbf k\cdot\mathbf x}
=-k^2e^{i\mathbf k\cdot\mathbf x}.
\]
Substitution into equation~\eqref{eq:c10-frw-wave} gives
\begin{equation}
\ddot\phi_{\mathbf k}
+3H\dot\phi_{\mathbf k}
+\frac{k^2}{a^2}\phi_{\mathbf k}=0.
\label{eq:c10-mode-equation}
\end{equation}
Thus every wave number behaves like a damped oscillator with
\begin{equation}
\gamma=3H,
\qquad
\omega_k(t)=\frac{k}{a(t)}.
\label{eq:c10-mode-identification}
\end{equation}
The physical wave number \(k/a\) redshifts as the universe expands. Consequently the effective restoring force \(k^2/a^2\) decreases toward zero.

A mode that begins at short physical wavelength has \(\omega_k\gg H\) and oscillates. Its wavelength is stretched, its frequency falls, and Hubble damping eventually dominates. It then approaches a nonoscillating value rather than being pulled back to zero. Different \(\mathbf k\) modes cross over at different times: a larger \(k\), corresponding at fixed \(a\) to a shorter wavelength, must wait for more expansion before its physical frequency falls to the Hubble scale.

\section{Horizon Exit and Frozen Vacuum Fluctuations}

Using the oscillator's critical condition as an order-one estimate,
\[
\gamma=2\omega_k,
\]
equation~\eqref{eq:c10-mode-identification} gives
\begin{equation}
3H=\frac{2k}{a},
\qquad
a_{\rm cross}=\frac{2k}{3H}.
\label{eq:c10-critical-mode}
\end{equation}
The numerical factors are not the conceptual point. The crossover occurs when the physical wave number \(k/a\) is of order \(H\).

Let
\[
\lambda_{\rm com}=\frac{2\pi}{k},
\qquad
\lambda_{\rm phys}=a\lambda_{\rm com}
=\frac{2\pi a}{k}.
\]
A dimensional check fixes the placement of \(a\): with dimensionless \(\mathbf x\), both \(k\) and \(\mathbf k\cdot\mathbf x\) are dimensionless, while \(a\) and \(\lambda_{\rm phys}\) carry length. Equation~\eqref{eq:c10-critical-mode} therefore implies
\begin{equation}
\lambda_{\rm phys}\sim H^{-1},
\label{eq:c10-horizon-crossing}
\end{equation}
up to factors such as \(2\), \(3\), and \(2\pi\).

The Hubble law \(v=Hd\) associates the distance \(H^{-1}\) with a recession speed of order \(c=1\). During inflation, \(H\) is large and nearly constant, so this Hubble length is microscopic even though it plays the role of a causal scale. Calling the event \emph{horizon exit} is conventional: no material object passes through a physical wall. The mode's physical wavelength grows beyond the approximately fixed Hubble scale, and its dynamics changes from rapid oscillation to an almost frozen amplitude.

\begin{figure}[htbp]
\centering
\begin{coursefitfigure}
\begin{tikzpicture}[x=1cm,y=0.85cm,>=Latex]
  \draw[->] (0,0) -- (9.2,0) node[right] {cosmic time};
  \draw[->] (0,0) -- (0,5.0) node[above] {physical length};
  \draw[dashed,thick] (0.25,2.75) -- (8.75,2.75)
    node[right] {\(H^{-1}\)};
  \draw[very thick,blue!55!black]
    (0.35,0.55) .. controls (3.3,0.75) and (5.3,1.55) .. (8.5,4.55);
  \node[blue!55!black,anchor=west] at (6.65,4.55) {\(\lambda_{\rm phys}=2\pi a/k\)};
  \draw[densely dotted] (6.05,0) -- (6.05,2.75);
  \fill (6.05,2.75) circle (2pt);
  \node[anchor=north] at (6.05,0) {horizon exit};
  \node[align=center] at (2.6,1.7) {oscillating\\sub-Hubble mode};
\end{tikzpicture}
\end{coursefitfigure}
\caption{During nearly exponential expansion, a physical wavelength grows while the Hubble length remains approximately fixed. Their crossing marks the transition from rapid oscillation to a frozen mode, up to order-one factors.}
\label{fig:c10-horizon-reconstruction}
\end{figure}

Vacuum fluctuations prevent inflation from simply erasing every wave. Every short-wavelength mode begins with zero-point motion. Expansion stretches one mode to the Hubble scale and freezes it, while still shorter vacuum modes continue to oscillate. They in turn stretch and freeze. Repetition builds a field profile containing frozen contributions over a broad range of wavelengths and orientations. The result is stochastic but not structureless: its statistics follow from the quantum state and the inflationary background.

\begin{figure}[htbp]
\centering
\includegraphics[width=0.88\textwidth]{lecture_10_figure_04.png}
\caption{A board-wide superposition of long and short field waves, with a smaller damping sketch beside the labels \(\gamma\) and \(3H\). The frame documents the qualitative mode picture; it is not used to read off an exact function.}
\label{fig:c10-mode-board}
\lectureframe{01:45:28.476}
\end{figure}

\begin{classroomqa}
\audiencequestion{Do the modes freeze at a point or at a wavelength, and is \(H\) itself growing exponentially?}
\lecturerresponse{They freeze when a wavelength reaches the Hubble scale. On the nearly flat plateau, \(V\) and hence \(H\) are approximately constant; it is the scale factor \(a(t)\), and therefore the physical wavelength \(a/k\), that grows exponentially. Each comoving mode reaches the condition \(k/a\sim H\) at its own time. \lecturetimestamp{01:45:06--01:45:39}}
\end{classroomqa}

\begin{classroomqa}
\audiencequestion{Are these electromagnetic waves?}
\lecturerresponse{No. They are fluctuations of the inflaton scalar field. Electromagnetic vacuum fluctuations provide a familiar analogy, but the variable in equation~\eqref{eq:c10-mode-equation} is \(\phi\). \lecturetimestamp{01:45:39--01:46:04}}
\end{classroomqa}

\begin{classroomqa}
\audiencequestion{Is there a continuum of wavelengths?}
\lecturerresponse{For the infinite spatial idealization, yes. The field is represented by a Fourier integral,
\begin{equation}
\phi(\mathbf x,t)
=\int\frac{d^3k}{(2\pi)^3}\,
\phi_{\mathbf k}(t)e^{i\mathbf k\cdot\mathbf x}.
\label{eq:c10-fourier-integral}
\end{equation}
Every wave number undergoes the same redshifting and freezing mechanism. \lecturetimestamp{01:46:04--01:46:34}}
\end{classroomqa}

\begin{classroomqa}
\audiencequestion{Would a finite universe have only discrete modes?}
\lecturerresponse{Yes. Boundary conditions on a finite spatial volume quantize the allowed wave numbers. For a volume much larger than the scales under discussion, the spacing can be so fine that the sum is well approximated by the integral in equation~\eqref{eq:c10-fourier-integral}. The continuous picture is an idealization, not a different local mechanism. \lecturetimestamp{01:46:34--01:47:01}}
\end{classroomqa}

Two motions now coexist. The spatially averaged field \(\bar\phi(t)\) rolls slowly and classically down the potential. On top of it sits a small quantum fluctuation,
\[
\phi(\mathbf x,t)=\bar\phi(t)+\delta\phi(\mathbf x,t),
\]
assembled from the frozen modes. The fluctuation energy itself is too small to be the observed matter contrast. Its importance is that it perturbs the time at which inflation ends.

A full single-field quantum calculation sharpens this oscillator argument. A light scalar acquires a typical fluctuation per logarithmic interval of wave number of order
\begin{equation}
\delta\phi_{\rm exit}\sim\frac{H}{2\pi}.
\label{eq:c10-field-amplitude}
\end{equation}
A derivation of the normalization requires quantizing the field modes in the inflationary background. The oscillator picture has already supplied the essential mechanism: vacuum fluctuations exist even in the ground state, and expansion converts them into a classical-looking frozen pattern.

\section{The Cliff Converts Field Noise into Density Noise}

First suppress \(\delta\phi\). At any fixed time the field then has one value everywhere, represented on the board by a vertical line when space \(x\) is drawn upward and field value \(\phi\) horizontally. As time advances, the line moves toward the cliff. The whole universe reaches the drop together, inflaton potential energy becomes particles, radiation, and heat, exponential expansion ends, and the resulting energy density begins to dilute with ordinary expansion.

\begin{figure}[htbp]
\centering
\includegraphics[width=0.88\textwidth]{lecture_10_figure_05.png}
\caption{The blackboard's space-versus-field diagram. The vertical axis is \(x\), the horizontal direction is \(\phi\), faint near-vertical profiles mark successive field configurations, and the hatched boundary at right is the cliff where inflation ends.}
\label{fig:c10-cliff-board}
\lectureframe{01:50:54.502}
\end{figure}

\begin{figure}[htbp]
\centering
\begin{coursefitfigure}
\begin{tikzpicture}[x=1cm,y=0.82cm,>=Latex]
  \draw[->] (0,0) -- (9.2,0) node[right] {\(\phi\)};
  \draw[->] (0,0) -- (0,5.2) node[above] {space \(x\)};
  \draw[thick,gray!60] (8.2,0.15) rectangle (8.55,4.85);
  \foreach \y in {0.35,0.85,...,4.55} {
    \draw[gray!60] (8.2,\y) -- (8.55,{\y+0.25});
  }
  \node[anchor=south,rotate=90] at (8.38,2.5) {cliff};
  \draw[thick,densely dashed,gray!70] (3.2,0.2) -- (3.2,4.75);
  \draw[thick,blue!45!black,smooth]
    plot coordinates {(5.0,0.2) (5.12,0.75) (4.9,1.25) (5.18,1.8) (4.94,2.35) (5.2,2.9) (4.88,3.45) (5.12,4.05) (5.0,4.75)};
  \draw[very thick,red!65!black,smooth]
    plot coordinates {(7.15,0.2) (7.32,0.7) (7.0,1.2) (7.42,1.75) (7.1,2.3) (7.5,2.85) (7.02,3.35) (7.35,4.0) (7.16,4.75)};
  \draw[->,thick] (3.55,2.5) -- (4.55,2.5);
  \draw[->,thick] (5.55,2.5) -- (6.65,2.5);
  \node[align=center,anchor=south,font=\small] at (3.2,4.8) {uniform\\profile};
  \node[align=center,anchor=south,font=\small] at (5.0,4.8) {frozen\\ripple};
  \node[align=center,anchor=south,font=\small] at (7.2,4.8) {near\\reheating};
\end{tikzpicture}
\end{coursefitfigure}
\caption{Clean reconstruction of the space-versus-field picture. The mean field moves toward the cliff while a frozen spatial ripple persists, so different positions reach the end of inflation at slightly different times.}
\label{fig:c10-cliff-reconstruction}
\end{figure}

Restore the frozen variation. The line is now slightly wavy: neighboring places have slightly different values of \(\phi\). Think of a rank of people trying to march shoulder to shoulder toward a cliff. A few are a little ahead, a few a little behind. The ones who are ahead go over first.

That timing offset is the real source of the density contrast. A region that goes over first reheats first. Its potential energy is dumped into particles and heat, and then that energy has time to dilute while the universe expands in the post-inflationary era. A nearby region that goes over slightly later reheats later, so when the two are compared, the later region is denser simply because it has had less time to dilute. The first set of crossings therefore leaves relatively dense pockets. While those pockets expand, shorter-wavelength ripples bring another set of positions over the edge and make smaller pockets inside the larger pattern. Repeating this across the full hierarchy of frozen modes produces structure within structure. Thus the tiny field fluctuation is not important because of its own energy; it is important because it shifts the local time at which reheating happens.

The timing relation can be written schematically. If the homogeneous field reaches the end value \(\phi_{\rm end}\) with velocity \(\dot{\bar\phi}\), then a small field displacement changes the local end time by
\begin{equation}
\delta t_{\rm end}
\simeq-\frac{\delta\phi}{\dot{\bar\phi}}.
\label{eq:c10-time-delay}
\end{equation}
A slower roll makes the same \(\delta\phi\) correspond to a larger delay. During that extra time the earlier region expands and dilutes more. Up to order-one transfer factors, the curvature or density perturbation is therefore controlled by
\begin{equation}
\mathcal R
\sim H\,\delta t_{\rm end}
\sim-\frac{H\,\delta\phi}{\dot{\bar\phi}}.
\label{eq:c10-curvature-estimate}
\end{equation}
Combining this with equation~\eqref{eq:c10-field-amplitude} gives the familiar slow-roll scale
\begin{equation}
\mathcal P_{\mathcal R}^{1/2}
\sim\frac{H^2}{2\pi|\dot{\bar\phi}|},
\label{eq:c10-slow-roll-amplitude}
\end{equation}
with precise conventions and subsequent transfer to \(\delta\rho/\rho\) requiring a fuller perturbation calculation. The slower the field rolls, the longer the time between neighboring cliff crossings and the greater the intervening dilution. The same field fluctuation therefore leaves a larger density imprint when \(\lvert\dot{\bar\phi}\rvert\) is smaller.

Many modes cross the Hubble scale at different times, so many differently sized regions reheat at slightly different times. The resulting density field contains overdense and underdense pockets nested across scales. Its statistics depend on the height, slope, and curvature of the inflationary potential. Suitable slow-roll models naturally produce a nearly scale-invariant amplitude near
\[
\frac{\delta\rho}{\rho}\sim10^{-5},
\]
and detailed CMB spectra allow a small number of cosmological parameters to be tested against many angular scales.

\section{From Primordial Seeds Back to the Sky}

We can now return to the observer at the center of a set of past-light shells. The last-scattering sphere cuts through the three-dimensional primordial density field. Directions in which that shell intersects different gravitational potentials produce slightly different microwave temperatures. WMAP displays the two-dimensional pattern on that sphere. The map is neither a photograph of time zero nor a projection of every structure inside the sphere; it is the temperature structure on the particular shell from which the microwave photons last scattered.

Inside the shell, the same primordial perturbations continue evolving. Gravity amplifies overdensities, caustic-like flows sharpen sheets and filaments, gas collects, and galaxies form. Thus the CMB pattern and the later cosmic web are two stages of one history: frozen inflationary fluctuations first appear as tiny density seeds and later become visible structure.

\begin{classroomqa}
\audiencequestion{What epoch does the microwave background we see today represent?}
\lecturerresponse{The photons have traveled for roughly \(13\) billion years, but they were released a few hundred thousand years after the hot Big Bang, when electrons and protons combined into neutral hydrogen and the universe became transparent. Inflation ended vastly earlier. Reheating created the hot matter and radiation; the universe then expanded and cooled until last scattering. For this argument the robust statement is the order of magnitude---a few hundred thousand years---rather than a precise recombination date. \lecturetimestamp{02:00:38--02:02:19}}
\end{classroomqa}

\section{One Sky and Cosmic Variance}

\begin{classroomqa}
\audiencequestion{Does inflation predict the reported preferred direction or north--south asymmetry in WMAP?}
\lecturerresponse{No such preferred direction follows from the simple inflationary picture developed here. The status of the reported large-angle anomalies was uncertain, foregrounds and statistical fluctuations were live concerns, and a responsible assessment required the specialist literature rather than a sensational summary. Inflation predicts statistical properties of an ensemble; deciding whether one unusual-looking sky feature is new physics is harder. \lecturetimestamp{02:02:19--02:03:26}}
\end{classroomqa}

Small angular features occur many times across the sky. They provide many samples, so their distribution and isotropy can be tested statistically. A quadrupole-sized feature is different: only a few independent large-angle modes fit on the observable sphere. Their measured power can therefore differ substantially from the ensemble mean even when the underlying theory is correct. This irreducible sampling uncertainty is \emph{cosmic variance}.

\begin{classroomqa}
\audiencequestion{What about the apparent alignment of several low multipoles?}
\lecturerresponse{It could be a clue, a foreground or analysis artifact, or an unlikely but allowed realization. The difficulty is that we cannot draw another observable universe from the same statistical ensemble. A coin giving four heads in six tosses need not be biased; repeated batches would decide, but cosmology supplies only one batch on the largest scales. The anomaly may deserve study, yet one sky cannot by itself provide the repeated trials needed for a decisive significance test. \lecturetimestamp{02:05:31--02:07:36}}
\end{classroomqa}

\section{From Quantum Motion to Cosmic Structure}

Three universal constants define the Planck scales and expose how far familiar particle physics lies from quantum gravity. The microwave sky supplies a precise cosmological target: near isotropy, a fluctuation amplitude around \(10^{-5}\), and structure spread over many scales. Quantum mechanics ensures that every field mode has zero-point fluctuations. Inflation stretches each inflaton mode until its physical wavelength is of order \(H^{-1}\), at which point Hubble damping suppresses its oscillation and leaves an effectively frozen amplitude.

The homogeneous inflaton continues to roll beneath that frozen noise. When the field reaches the end of its plateau, neighboring regions do not reheat at exactly the same time. Earlier regions dilute longer; later regions remain denser. This converts \(\delta\phi\) into a primordial density and curvature pattern. Gravity then amplifies the pattern into galaxies and the cosmic web, while the last-scattering shell preserves its earlier state in the microwave sky. Microscopic zero-point motion, cosmic expansion, and staggered reheating are therefore not separate stories. Together they explain how an inflationary universe can be both extraordinarily smooth and structured enough for us to exist.
